\subsection{Bestehende Plattformen im Vergleich}\label{sec:wettbewerbsanalyse}

Chefkoch.de zählt zu den etablierten deutschsprachigen Rezept-Communities für Hobbyköch:innen. Nutzer:innen laden dort eigene Rezepte hoch, bewerten fremde Beiträge und kommentieren sie; dieser tatsächliche Austausch begründet die Einstufung \enquote{stark} bei Kommunikation \& Community in~\autoref{tab:wettbewerbsvergleich}. Für Resteverwertung bietet die Plattform zwar Rezeptsammlungen zu einzelnen Zutaten, aber kein Werkzeug, das aus mehreren vorhandenen Resten Vorschläge ableitet.

Auf TikTok, das hier stellvertretend für Instagram und ähnliche Plattformen steht, hat sich Kochen zu einer eigenen Content-Kategorie mit hoher Reichweite kurzer Rezeptvideos entwickelt. Kommentare, geteilte Videos und das Folgen einzelner Kanäle sorgen für intensive Interaktion; ein Werkzeug für die gezielte Suche nach vorhandenen Zutaten fehlt jedoch ebenso wie eine eigene Reste-Funktion.

Eine zweite Plattform-Kategorie bilden spezialisierte Reste-Matching-Apps wie Restegourmet und die \ac{BMEL}-App \enquote{Zu gut für die Tonne!}. Beide nehmen mehrere übrige Zutaten entgegen und schlagen passende Rezepte vor; die \ac{BMEL}-App ergänzt das um einen Portionenrechner und eine umfangreiche Rezeptdatenbank. Eine echte Social-Community bieten sie aber nicht: Following, Kommentar-Communities und Gamification fehlen, und auch Restegourmets Teilen und Bewerten einzelner Rezepte schafft keinen Ort für Austausch.

Die dritte Reste-Matching-App, leftovercooking, kommt Resterias Ansatz am nächsten: Sie verbindet Zutaten-Matching bereits mit ersten Community-Elementen, erfasst den vorhandenen Lebensmittelvorrat, markiert bald verderbliche Zutaten farblich und erzeugt daraus KI-gestützte Rezeptvorschläge. Mehrere Food-Creator:innen teilen dort eigene Rezepte, und ein Belohnungspunkte-System zeigt die durch verwertete Reste gesparten Geldbeträge an – individuelles Spar-Feedback ohne Bezug zu anderen Nutzer:innen, das nicht als die im Kriterium gemeinte, sozial-interaktive Gamification zählt. Ein Austausch zwischen gewöhnlichen Nutzer:innen, etwa über Kommentare oder gegenseitiges Folgen, ist nicht belegt; die Community-Ebene bleibt auf das Verfolgen von Creator-Inhalten beschränkt.

Foodsharing.de scheint Community und Anti-Verschwendungs-Fokus zu verbinden, organisiert aber nur die Weitergabe überschüssiger Lebensmittel über Fairteiler (öffentliche Abgabestellen) und direkte Abholungen, ohne aus Resten Rezepte vorzuschlagen. Die Mitglieder-Community dort ist aktiv und real, was die Einstufung \enquote{stark} rechtfertigt, bezieht sich aber auf die Weitergabe vor dem Verderb, nicht auf die Verarbeitung in der eigenen Küche.

Grundlage der Einstufungen in~\autoref{tab:wettbewerbsvergleich} ist eine eigene Sichtung der Web-Auftritte und Apps aller sechs Plattformen im Juli 2026. Die Kriterien leiten sich aus den in der Aufgabenstellung genannten Achsen Rezept-Austausch, -Entdeckung und Kommunikation ab, ergänzt um Nachhaltigkeits-/Restefokus und Reste-Matching-Werkzeug. \enquote{Stark} bedeutet ein aktiv genutztes, zentrales Merkmal, \enquote{mittel} ein vorhandenes, aber eingeschränktes und \enquote{schwach} ein kaum oder nicht vorhandenes. Bei Kommunikation \& Community zählt nur tatsächlicher Austausch zwischen Nutzer:innen; ein bloßes Bewertungsfeld reicht dafür nicht aus. Rezept-Entdeckung erfasst, wie leicht sich überhaupt ein passendes Rezept finden lässt, während die reste-spezifische Suche das Kriterium Reste-Matching-Werkzeug abdeckt.

\begin{table}[H]
\caption{Vergleich bestehender Plattformen anhand zentraler Kriterien}\label{tab:wettbewerbsvergleich}
\begin{tabularx}{\textwidth}{@{}l>{\centering\arraybackslash}X>{\centering\arraybackslash}X>{\centering\arraybackslash}X>{\centering\arraybackslash}X>{\centering\arraybackslash}X@{}}
\toprule
\textbf{Plattform} & \textbf{Rezept-Austausch} & \textbf{Rezept-Entdeckung} & \textbf{Kommunikation \& Community} & \textbf{Nach\-hal\-tig\-keits-/\allowbreak Restefokus} & \textbf{Reste-Matching-Werkzeug} \\
\midrule
Chefkoch.de & stark & stark & stark & mittel & schwach \\
TikTok (Food) & mittel & stark & stark & schwach & schwach \\
Restegourmet & mittel & schwach & schwach & stark & stark \\
\ac{BMEL}-App & mittel & stark & schwach & stark & stark \\
leftovercooking & mittel & mittel & mittel & stark & stark \\
Foodsharing.de & schwach & schwach & stark & stark & schwach \\
\bottomrule
\end{tabularx}
\quelle{Eigene Darstellung auf Basis eigener Sichtung der Web-Auftritte und Apps von Chefkoch, TikTok, Restegourmet, Zu gut für die Tonne, leftovercooking und Foodsharing, Juli 2026.}
\end{table}

Die Tabelle macht die Lücke sichtbar, die \enquote{Resteria} schließen soll: Die Matching-Apps lösen das Zutaten-Problem ohne Austausch zwischen den Nutzer:innen, die Community-Plattformen den Austausch ohne Weg von übrigen Zutaten zum Rezept. Die Stärke der eigenen Konzeptposition liegt in dieser Kombination, die Schwäche im fehlenden Nutzertest; Chancen bieten die wachsende Relevanz von Lebensmittelverschwendung und Partnerschaften mit bestehenden Initiativen, ein Risiko das dichte Umfeld etablierter Reste-Apps.

Der gewichtigste Einwand gegen diese Einschätzung lautet, eine reine Community-Plattform mit Nachhaltigkeitsthema könnte bereits ausreichen. Ohne Matching-Werkzeug müssten Nutzer:innen dort jedoch weiterhin selbst nach passenden Rezepten suchen, und genau diese Hürde nehmen die Reste-Matching-Apps ihnen ab. Das Alleinstellungsmerkmal von Resteria liegt damit in der Verbindung beider Elemente um eine wechselseitige Community herum, die über das Verfolgen von Creator-Inhalten hinausgeht.
